\chapter{A0 序 感谢 前言} % Introduction chapter suppressed from the table of contents

\hypertarget{ux5e8f-preface}{%
\section{序 Preface}\label{ux5e8f-preface}}

\hypertarget{ux524dux8a00-prologue}{%
\section{前言 Prologue}\label{ux524dux8a00-prologue}}

"There are no hard distinctions between what is real and what is unreal,
nor between what is true and what is false. A thing is not necessarily
either true or false; it can be both true and false."

\begin{description}
\item[]
\begin{description}
\tightlist
\item[]
Harold Pinter , 1958
\end{description}
\end{description}

我是从自学社会系统理论 (socio-technical systems theory) 开始听到 Russell
Ackoff 的名字。
开始关注他的著作是从看完他的演讲视频开始：觉得他演讲很清晰，针对主题用实例解释，而且观点都很精炼，并且不落书套，已经经过深思熟虑的结果。我便开始搜索他简历与著作：

原来在70至80年代，他在宾州大学 (U. of Pennsylvania) 沃顿 (Wharton
School) 举办 Social Systems Science Program
课程，帮助学生学习与研究社会系统。这课程被后人评价：``能集合理论与实际，帮助学员能把学到的用于自己熟识的领域。''

他出版了超过40本书，除了第一本《运筹学》(Operations
Research，1957)外，其他大部分都与社会系统或管理相关。

2007年，他 Herbert J. Addison 与 Sally Bibb 出版 "Management f-Laws :
How organizations really
work"：（也是他的最后一本书）。当我在图书馆里找到这本很薄的小书，发现与他之前出版的书很不同：

\begin{itemize}
\tightlist
\item
  风格轻松，之前的书，都是偏学术性，以学者的口吻写，这小书就显得幽默
\item
  简单易读：里面简单每2页介绍一条，共81条 f-laws.
\end{itemize}

他解释说： ``f-laws are truths about organizations that we might wish to
deny or ignore -- simple and more reliable guides to managers' everyday
behavior than the complex truths proposed by scientists, economists,
sociologists, politicians and philosophers.''

但这小书也有不足：每一条右边一页都是 Sally Bibb
的反馈意见，但都是基于她自己的想法，缺乏实例，有些更只是用另一个角度把左边A的点重新说。

所以我便基于古人写《左传》（注）的方式，依据 Ackoff
教授每条经典道理，加上我自己的经历或咨询过企业的故事，希望能帮助读者能更好理解他每条的重点。

我写《敏捷开发A++
实践》时，曾经收到反馈说有些章太长（十几页），如能压缩可以提升易读性，所以在这小册里，把每条的相关故事都控制在一页内。

基于Ackoff
教授的XX条基础，我后面继续尝试其他管理大师，如XXX，也写出f-laws，所以这小册取名为《管理
A++》。

正如英国剧本作家Harold
Pinter说，任何事情没有绝对的对或错，因为各企业或团队的环境、能力不同，这小册里的道理不一定都适用。只希望站在巨人的肩膀上，借用实例，让读者更了解道理背后的原意，并得到启发。

\begin{description}
\tightlist
\item[]
注：《春秋》是孔子根據魯國史官所編之史書重新修訂而成，記述從魯隱公元年（公元前722年）到魯哀公十四年（公元前481年）間二百四十二年之歷史。但《春秋》文字过于简质，后人不易理解，所以相继各种解释和说明，称之为``传''，例如《公羊传》，其中《左传》主要以写故事的方式来解读《春秋》，
\end{description}

\hypertarget{ux53c2ux8003-references}{%
\section{参考 References}\label{ux53c2ux8003-references}}

\begin{enumerate}
\tightlist
\item
  Ackoff, Russell, et al. \emph{Management f-Laws : How organizations
  really work}(Triarchy Press, 2007)\\
\end{enumerate}


